\section*{\translate{Terminologi}}\label{sec:terminology}
\phantomsection
\addcontentsline{toc}{section}{Terminologi}

\begin{table}[ht!]
  \renewcommand{\arraystretch}{1.8}
  \begin{tabular}{l p{10cm}}
  \toprule
  \textbf{Term} & \textbf{Förklaring} \\
  \midrule
  Agenter & Ett arkitekturmönster där LLM:en är utrustad med verktyg som den kan använda för att utföra uppgifter. \\
  API & \emph{Application Programming Interface} En uppsättning regler och protokoll som tillåter olika programvaror att kommunicera med varandra.\\
  Design Science Methology & En forskningsmetod som fokuserar på att skapa artefakter som utvärderas för att lösa och identifiera problem.\\
  HTTP & \emph{Hypertext Transfer Protocol} Ett protokoll som används för att överföra data över webben.\\

  LLM & \emph{Large Language Model} En typ av AI-modell tränad på stora mängder text som kan generera och förstå naturligt språk.\\
  MCP & \emph{Model Context Protocol} Ett öppet standardprotokoll utvecklat av Anthropic som möjliggör för LLM-agenter att kommunicera med externa verktyg och datakällor på ett enhetligt sätt.\\
  PoC & \emph{Proof Of Concept} En demonstration som visar att en viss idé eller koncept är genomförbart.\\
  SDK & \emph{Software Development Kit} En samling verktyg, bibliotek och dokumentation som underlättar utvecklingen av programvara för en specifik plattform eller ramverk.\\
  TypeScript & Ett programmeringsspråk \\
  Verktyg & En funktion eller resurs som kan användas av en LLM för att utföra en specifik uppgift, såsom att hämta data från en databas eller utföra en beräkning.\\
  \bottomrule
  \end{tabular}
\end{table}

